\section{Formal Theory for CVaR-CJE}
\label{app:formal}
\label{app:proofs}

This appendix gives a self-contained definition of mean-CJE and CVaR-CJE.
It starts from the probability laws induced by a policy, then derives the
mean and lower-tail objects, the direct shortfall estimator, and the audit
moments used to decide whether a calibrated level claim is justified.

\subsection{Notation and Probability Laws}

Let $P_X$ be the prompt law on a measurable space $\mathcal{X}$. A policy
$\ptarget$ is a Markov kernel from prompts to answers. Under policy $\ptarget$,
\[
\X\sim P_X,\qquad
\A_\ptarget\sim \ptarget(\cdot\mid \X),\qquad
\Sscore_\ptarget=s(\X,\A_\ptarget),\qquad
\Y_\ptarget=y(\X,\A_\ptarget)\in[0,1].
\]
The variable $\Sscore$ is the cheap judge score and $\Y$ is the oracle outcome.
Write
\[
O_\ptarget=(\X,\A_\ptarget,\Sscore_\ptarget,\Y_\ptarget)
\]
and let $P_\ptarget$ be the joint law of $O_\ptarget$. Expectations and
probabilities under $P_\ptarget$ are denoted $\E_\ptarget$ and $\Pp_\ptarget$.
The reference policy is $\pzero$ and the target policy is $\ptarget$.

Let
\[
\mathcal{G}=\sigma(\Sscore,\X)
\]
be the sigma-field generated by the prompt and cheap score. The reference-law
conditional mean and reference-law conditional shortfall are
\begin{align}
f_0(\Sscore,\X)
&=\E_{\pzero}[\Y\mid \mathcal{G}],\\
h_{0,t}(\Sscore,\X)
&=\E_{\pzero}\!\left[\shortfall{t}{\Y}\mid \mathcal{G}\right],
\qquad t\in\R.
\end{align}
When the reference law is clear, write $f=f_0$ and $h_t=h_{0,t}$. The mean
residual and shortfall residual are
\begin{align}
\varepsilon_f
&=\Y-f(\Sscore,\X),\\
u_t
&=\shortfall{t}{\Y}-h_t(\Sscore,\X).
\end{align}
Conditional equality under $\mathcal{G}$ is sufficient for transport, but the
identification statements below only require target-law mean residuals to
vanish.

\subsection{Mean-CJE}

The mean value of policy $\ptarget$ is
\begin{equation}
V(\ptarget)=\E_\ptarget[\Y].
\end{equation}
For any measurable surrogate $f(\Sscore,\X)$, define the surrogate value
\begin{equation}
V_f(\ptarget)=\E_\ptarget[f(\Sscore,\X)].
\end{equation}
Then
\begin{equation}
V(\ptarget)-V_f(\ptarget)
=\E_\ptarget[\Y-f(\Sscore,\X)]
=\E_\ptarget[\varepsilon_f].
\label{eq:mean-transport-identity}
\end{equation}

\begin{assumption}[Mean transport]
\label{ass:mean-transport}
For the target policy $\ptarget$,
\begin{equation}
\E_\ptarget[\Y-f(\Sscore,\X)]=0.
\end{equation}
\end{assumption}

\begin{proposition}[Mean-CJE identification]
\label{prop:formal-mean-identification}
For any fixed surrogate $f$, \cref{ass:mean-transport} holds if and only if
$V_f(\ptarget)=V(\ptarget)$.
\end{proposition}

\begin{twocolproof}
\twocolrow{
\[
V(\ptarget)-V_f(\ptarget)
=\E_\ptarget[\Y]-\E_\ptarget[f(\Sscore,\X)].
\]
}{
Use the definitions of the oracle value and surrogate value.
}
\twocolrow{
\[
\E_\ptarget[\Y]-\E_\ptarget[f(\Sscore,\X)]
=\E_\ptarget[\Y-f(\Sscore,\X)].
\]
}{
Linearity of expectation applies because both terms are integrable.
}
\twocolrow{
\[
V_f(\ptarget)=V(\ptarget)
\quad\Longleftrightarrow\quad
\E_\ptarget[\Y-f(\Sscore,\X)]=0.
\]
}{
Move the previous identity to the zero-difference form. This is exactly the
mean-transport condition.
}
\end{twocolproof}

The corresponding audit target is the scalar moment
\[
m_f(\ptarget)=\E_\ptarget[\varepsilon_f].
\]
A mean-CJE level claim is justified for $\ptarget$ only when this target-policy
moment is indistinguishable from zero at the desired level of precision.

\subsection{Lower-Tail CVaR}

Fix $\alpha\in(0,1]$. Let
\[
F_\ptarget(t)=\Pp_\ptarget(\Y\le t),
\qquad
q_\alpha(\ptarget)=\inf\{t:F_\ptarget(t)\ge \alpha\}.
\]
The lower-tail stop-loss objective is
\begin{equation}
\Psi_{\alpha,\ptarget}(t)
=t-\frac{1}{\alpha}\E_\ptarget\!\left[\shortfall{t}{\Y}\right].
\label{eq:app-stoploss-cvar}
\end{equation}
The lower-tail CVaR is
\begin{equation}
\CVaR_\alpha(\ptarget)=\sup_{t\in\R}\Psi_{\alpha,\ptarget}(t).
\label{eq:app-cvar-definition}
\end{equation}

\begin{assumption}[Regular lower tail]
\label{ass:regular-tail}
For $\alpha\in(0,1)$, the target distribution of $\Y$ is continuous at
$q_\alpha(\ptarget)$, has positive density in a neighborhood of
$q_\alpha(\ptarget)$, and the maximizer of \cref{eq:app-stoploss-cvar} is
unique up to quantile equivalence.
\end{assumption}

\begin{proposition}[Stop-loss derivation]
\label{prop:formal-stoploss}
Under \cref{ass:regular-tail}, the optimizer of
\cref{eq:app-stoploss-cvar} is $t_\alpha^\star=q_\alpha(\ptarget)$, and
\begin{equation}
\CVaR_\alpha(\ptarget)
=\E_\ptarget[\Y\mid \Y\le q_\alpha(\ptarget)].
\end{equation}
\end{proposition}

\begin{twocolproof}
\twocolrow{
\[
\frac{d}{dt}\E_\ptarget[\shortfall{t}{\Y}]
=\Pp_\ptarget(\Y\le t)=F_\ptarget(t).
\]
}{
At continuity points, differentiating $(t-\Y)_+$ in $t$ gives
$\one\{\Y\le t\}$, and \cref{ass:regular-tail} permits interchange with
expectation.
}
\twocolrow{
\[
\frac{d}{dt}\Psi_{\alpha,\ptarget}(t)
=1-\frac{F_\ptarget(t)}{\alpha}.
\]
}{
Differentiate the stop-loss objective in \cref{eq:app-stoploss-cvar}.
}
\twocolrow{
$\Psi_{\alpha,\ptarget}$ increases when $F_\ptarget(t)<\alpha$ and decreases
when $F_\ptarget(t)>\alpha$.
}{
The sign of the derivative is determined by whether the target tail mass has
not yet reached, or has already exceeded, $\alpha$.
}
\twocolrow{
\[
t_\alpha^\star=q_\alpha(\ptarget).
\]
}{
Under regularity, the sign change is unique up to quantile equivalence.
}
\twocolrow{
\[
\begin{aligned}
\Psi_{\alpha,\ptarget}(q_\alpha)
&=q_\alpha-\frac{1}{\alpha}
\E_\ptarget[(q_\alpha-\Y)\one\{\Y\le q_\alpha\}]\\
&=\frac{1}{\alpha}\E_\ptarget[\Y\one\{\Y\le q_\alpha\}].
\end{aligned}
\]
}{
Continuity gives $\Pp_\ptarget(\Y\le q_\alpha)=\alpha$, so the
$q_\alpha$ terms cancel after expanding the shortfall.
}
\twocolrow{
\[
\Psi_{\alpha,\ptarget}(q_\alpha)
=\E_\ptarget[\Y\mid \Y\le q_\alpha].
\]
}{
Divide the truncated first moment by the tail probability $\alpha$.
}
\end{twocolproof}

When $\Y$ has atoms, the objective remains the primitive. The subdifferential
of $t\mapsto \E_\ptarget[\shortfall{t}{\Y}]$ at $t$ is
\[
\left[\Pp_\ptarget(\Y<t),\Pp_\ptarget(\Y\le t)\right],
\]
so $t$ is optimal exactly when
\[
\Pp_\ptarget(\Y<t)\le \alpha \le \Pp_\ptarget(\Y\le t).
\]
The optimized value is the partial-mass expected shortfall derived in
\cref{app:discrete}.

\subsection{Discrete, Count, and Nonnegative Outcomes}

The stop-loss definition is the correct primitive for non-continuous laws. The
general-loss treatment of \citet{RockafellarUryasev2002} emphasizes that, when
the distribution has a mass point at the VaR threshold, CVaR is not the naive
conditional mean above or below that threshold. The boundary atom must be split
so that exactly the requested tail probability is averaged. This is also the
reason the optimization formula is preferable to a verbal definition such as
``the mean beyond VaR'' \citep{RockafellarUryasev2000,AcerbiTasche2002}.

For an upper-tail loss $L$ where larger values are worse, let
$F_L(t)=\Pp(L\le t)$, fix $\beta\in(0,1)$, and set
\[
q_\beta^L=\inf\{t:F_L(t)\ge\beta\}.
\]
The atom-split upper-tail expected shortfall is
\begin{equation}
\ES_\beta^{\mathrm{upper}}(L)
=\frac{1}{1-\beta}
\left(
\E[L\one\{L>q_\beta^L\}]
+q_\beta^L\left(F_L(q_\beta^L)-\beta\right)
\right).
\label{eq:upper-discrete-es}
\end{equation}
For a lower-tail score $\Y$ where smaller values are worse, the corresponding
left-tail formula is \cref{eq:discrete-es}. The two conventions are equivalent
after applying the upper-tail loss formula to $L=-\Y$.

For integer-valued or count outcomes, the VaR threshold is integer-valued, but
the expected shortfall is generally real-valued because only part of the
boundary atom is used. This distinction matters for Poisson, negative-binomial,
zero-inflated, and bounded-count laws. In count-process forecasting,
\citet{HomburgWeissAlwanFrahmGob2021} distinguish risk forecasts that compute
the risk functional from the discrete conditional PMF from Gaussian
approximations to the count process; their simulations find that Gaussian
approximations can materially distort tail-risk forecasts, especially through
risk underrating. For CVaR-CJE this supports a simple rule: if oracle outcomes
are integer-valued or rounded, use the discrete PMF or empirical mass function
directly, and do not smooth away the atom at the selected threshold merely to
recover the continuous formula.

For nonnegative outcomes that are not counts, exponential-dispersion families
provide a useful modeling bridge. \citet{LandsmanValdez2005} derive tail
conditional expectation formulas for exponential-dispersion models, including
continuous nonnegative laws such as Gamma and inverse Gaussian, and discrete
laws such as Poisson, binomial, and negative binomial. These results are not
needed for identification, but they clarify how a parametric sensitivity
analysis could be built when the oracle outcome is a nonnegative loss rather
than a bounded score.

Finally, the optimization literature gives two computational lessons. First,
sample-based CVaR and CVaR constraints can be written as convex programs, often
linear programs after adding auxiliary variables
\citep{RockafellarUryasev2000,RockafellarUryasev2002}. Second, CVaR regression
can be formulated through Rockafellar error and mixed-quantile quadrangles,
giving a linear-programming route to estimate CVaR as a function of covariates
\citep{GolodnikovKuzmenkoUryasev2019}. Recent large-scale work on
CVaR-constrained quadratic programs likewise treats the sample CVaR constraint
as a convex object and focuses on fast projection algorithms
\citep{LuxenbergPerezPineiroDiamondBoyd2025}. In the present paper these
computational facts justify the use of threshold-indexed shortfall calibration;
the statistical audit remains separate from the convex representation.

\subsection{Direct CVaR-CJE Population Target}

For each threshold $t$, define the shortfall variable
\begin{equation}
Z_t=\shortfall{t}{\Y}.
\end{equation}
CVaR estimation is therefore a continuum of mean estimation problems indexed
by $t$. The reference shortfall surrogate is
\begin{equation}
h_t(\Sscore,\X)=\E_{\pzero}[Z_t\mid \mathcal{G}].
\end{equation}
For a target policy $\ptarget$, the surrogate stop-loss objective is
\begin{equation}
\Psi^h_{\alpha,\ptarget}(t)
=t-\frac{1}{\alpha}\E_\ptarget[h_t(\Sscore,\X)].
\label{eq:formal-surrogate-stoploss}
\end{equation}

\begin{assumption}[Shortfall transport]
\label{ass:shortfall-transport}
On an index set $\mathcal{T}\subseteq\R$, the target policy satisfies
\begin{equation}
\E_\ptarget\!\left[\shortfall{t}{\Y}-h_t(\Sscore,\X)\right]=0
\qquad\text{for every }t\in\mathcal{T}.
\end{equation}
\end{assumption}

\begin{theorem}[CVaR-CJE identification]
\label{thm:formal-cvar-identification}
If \cref{ass:shortfall-transport} holds on $\mathcal{T}$, then
\[
\Psi^h_{\alpha,\ptarget}(t)=\Psi_{\alpha,\ptarget}(t)
\qquad\text{for every }t\in\mathcal{T}.
\]
Consequently,
\[
\sup_{t\in\mathcal{T}}\Psi^h_{\alpha,\ptarget}(t)
=
\sup_{t\in\mathcal{T}}\Psi_{\alpha,\ptarget}(t).
\]
If $\mathcal{T}=\R$, the direct CVaR-CJE population target equals
$\CVaR_\alpha(\ptarget)$. If $\mathcal{T}$ is a neighborhood containing the
oracle maximizers and the maximization is restricted to $\mathcal{T}$, the
same equality holds locally. Under \cref{ass:regular-tail}, the common local
maximizer is $t_\alpha^\star=q_\alpha(\ptarget)$.
\end{theorem}

\begin{twocolproof}
\twocolrow{
For any $t\in\mathcal{T}$,
\[
\begin{aligned}
\Psi_{\alpha,\ptarget}(t)-\Psi^h_{\alpha,\ptarget}(t)
&=-\frac{1}{\alpha}
\E_\ptarget[\shortfall{t}{\Y}-h_t(\Sscore,\X)].
\end{aligned}
\]
}{
Subtract the surrogate stop-loss objective from the oracle stop-loss objective;
the threshold terms cancel.
}
\twocolrow{
\[
\Psi_{\alpha,\ptarget}(t)-\Psi^h_{\alpha,\ptarget}(t)=0.
\]
}{
Apply \cref{ass:shortfall-transport} at the same threshold $t$.
}
\twocolrow{
\[
\sup_{t\in\mathcal{T}}\Psi^h_{\alpha,\ptarget}(t)
=\sup_{t\in\mathcal{T}}\Psi_{\alpha,\ptarget}(t).
\]
}{
Identical functions on the same index set have identical suprema on that set.
}
\twocolrow{
If $\mathcal{T}=\R$, the common value is
$\CVaR_\alpha(\ptarget)$.
}{
Use the definition of CVaR as the global supremum of the oracle stop-loss
objective.
}
\twocolrow{
Under \cref{ass:regular-tail}, the common local optimizer is
$q_\alpha(\ptarget)$.
}{
The stop-loss derivation in \cref{prop:formal-stoploss} identifies the regular
optimizer.
}
\end{twocolproof}

Shortfall transport is stronger than mean transport because it requires a
threshold-indexed family of residual moments to vanish. Mean-CJE uses only
$\E_\ptarget[\Y-f(\Sscore,\X)]=0$.

\subsection{Finite-Sample Direct Estimator}

Let $\mathcal{T}_n$ be a fixed finite threshold grid. For each
$t\in\mathcal{T}_n$, let $\hat h_t$ be a fitted shortfall calibrator. On an
independent target-policy evaluation sample
$(\X_i,\A_i,\Sscore_i)_{i=1}^n$, define
\begin{equation}
\hat\Psi_{\alpha,\ptarget}(t)
=t-\frac{1}{\alpha n}\sum_{i=1}^n \hat h_t(\Sscore_i,\X_i).
\label{eq:formal-plugin-objective}
\end{equation}
The selected threshold and direct CVaR-CJE estimate are
\begin{align}
\hat t_\alpha
&\in \operatorname*{arg\,max}_{t\in\mathcal{T}_n}
\hat\Psi_{\alpha,\ptarget}(t),\\
\widehat{\CVaR}_{\alpha}^{\mathrm{direct}}(\ptarget)
&=\hat\Psi_{\alpha,\ptarget}(\hat t_\alpha).
\end{align}

\begin{assumption}[Nuisance and grid regularity]
\label{ass:nuisance}
The shortfall calibrators are fitted on folds independent of the evaluation
folds, are uniformly square-integrable on neighborhoods of the relevant
maximizers, and satisfy
\[
\sup_{t\in\mathcal{T}_n}
\left|
\frac{1}{n}\sum_{i=1}^n \hat h_t(\Sscore_i,\X_i)
-\E_\ptarget[h_t(\Sscore,\X)]
\right|
=o_p(1).
\]
The grid $\mathcal{T}_n$ is fixed before policy comparisons and becomes dense
enough around the population maximizers that grid approximation error is
negligible at the target precision.
\end{assumption}

Under \cref{ass:nuisance}, $\hat\Psi_{\alpha,\ptarget}(t)$ uniformly
approximates the surrogate objective on the grid. Under
\cref{ass:shortfall-transport}, the surrogate objective is the oracle
stop-loss objective on the relevant index set, so the plug-in maximizer targets
the lower-tail oracle value.

\subsection{Audit Moments and Wald Statistic}
\label{app:audit}

The audit checks two population restrictions at the selected threshold:
tail mass and shortfall transport. At a regular population optimizer
$t_\alpha^\star$, define
\begin{align}
g_1(O_\ptarget;t_\alpha^\star)
&=\one\{\Y\le t_\alpha^\star\}-\alpha,\\
g_2(O_\ptarget;t_\alpha^\star)
&=\shortfall{t_\alpha^\star}{\Y}
-h_{t_\alpha^\star}(\Sscore,\X).
\end{align}

\begin{theorem}[Two-moment audit target]
\label{thm:formal-two-moment}
Under \cref{ass:regular-tail,ass:shortfall-transport},
\[
\E_\ptarget[g_1(O_\ptarget;t_\alpha^\star)]=0,
\qquad
\E_\ptarget[g_2(O_\ptarget;t_\alpha^\star)]=0.
\]
The first moment is the quantile restriction, and the second moment is the
shortfall-transport restriction at the selected threshold.
\end{theorem}

\begin{twocolproof}
\twocolrow{
\[
t_\alpha^\star=q_\alpha(\ptarget),
\qquad
\Pp_\ptarget(\Y\le t_\alpha^\star)=\alpha.
\]
}{
\cref{ass:regular-tail} gives a regular lower-tail quantile optimizer.
}
\twocolrow{
\[
\E_\ptarget[g_1(O_\ptarget;t_\alpha^\star)]
=\Pp_\ptarget(\Y\le t_\alpha^\star)-\alpha=0.
\]
}{
Take expectation of the indicator moment and substitute the target tail mass.
}
\twocolrow{
\[
\E_\ptarget[g_2(O_\ptarget;t_\alpha^\star)]
=\E_\ptarget[
\shortfall{t_\alpha^\star}{\Y}
-h_{t_\alpha^\star}(\Sscore,\X)].
\]
}{
This is the definition of the shortfall residual at the selected threshold.
}
\twocolrow{
\[
\E_\ptarget[g_2(O_\ptarget;t_\alpha^\star)]=0.
\]
}{
If $t_\alpha^\star\in\mathcal{T}$, this is exactly
\cref{ass:shortfall-transport} evaluated at $t_\alpha^\star$.
}
\end{twocolproof}

In a held-out audit sample with oracle outcomes, use the fitted threshold and
calibrator to form
\begin{align}
\hat g_{1i}
&=\one\{\Y_i\le \hat t_\alpha\}-\alpha,\\
\hat g_{2i}
&=\shortfall{\hat t_\alpha}{\Y_i}
-\hat h_{\hat t_\alpha}(\Sscore_i,\X_i),
\end{align}
and let
\[
\bar g_n=\frac{1}{n}\sum_{i=1}^n
\begin{pmatrix}
\hat g_{1i}\\
\hat g_{2i}
\end{pmatrix}.
\]
Let $\Omega$ be the asymptotic covariance of $\sqrt{n}\bar g_n$, including the
first-order effect of fitting $\hat h_t$ and selecting $\hat t_\alpha$, and let
$\hat\Omega$ be a consistent estimator of $\Omega$. The Wald statistic is
\begin{equation}
W_n=n\bar g_n^\top\hat\Omega^{-1}\bar g_n.
\label{eq:formal-wald}
\end{equation}

\begin{theorem}[Wald limit]
\label{thm:formal-wald-limit}
Suppose the two audit moments have mean zero, \cref{ass:nuisance} holds, and
the joint moment expansion satisfies
\[
\sqrt{n}\bar g_n
=\frac{1}{\sqrt{n}}\sum_{i=1}^n \phi(O_i)+o_p(1),
\qquad
\E[\phi(O_i)]=0,\qquad
\E[\phi(O_i)\phi(O_i)^\top]=\Omega,
\]
with $\Omega$ nonsingular and $\hat\Omega\overset{p}{\to}\Omega$. Then
\[
W_n\rightsquigarrow \chi^2_2.
\]
\end{theorem}

\begin{twocolproof}
\twocolrow{
\[
\sqrt{n}\bar g_n
=\frac{1}{\sqrt{n}}\sum_{i=1}^n\phi(O_i)+o_p(1).
\]
}{
This is the assumed first-order expansion of the estimated audit moments.
}
\twocolrow{
\[
\sqrt{n}\bar g_n\rightsquigarrow N(0,\Omega).
\]
}{
Apply the multivariate central limit theorem to the influence-function sum.
}
\twocolrow{
\[
\hat\Omega^{-1}\overset{p}{\to}\Omega^{-1}.
\]
}{
Covariance consistency and nonsingularity allow inversion by the continuous
mapping theorem.
}
\twocolrow{
\[
n\bar g_n^\top\hat\Omega^{-1}\bar g_n
\rightsquigarrow Z^\top\Omega^{-1}Z,\quad Z\sim N(0,\Omega).
\]
}{
Use Slutsky's theorem and the continuous mapping theorem on the quadratic form.
}
\twocolrow{
\[
Z^\top\Omega^{-1}Z=U^\top U\sim\chi^2_2,
\qquad U\sim N(0,I_2).
\]
}{
Whiten $Z$ as $Z=\Omega^{1/2}U$. The two components of $U$ correspond to the
two audit moments.
}
\end{twocolproof}

At audit level $\eta$, reject the calibrated level claim when
\[
W_n>\chi^2_{2,1-\eta}.
\]
A rejection means the target policy has failed at least one of the two required
restrictions: the threshold does not have the intended lower-tail mass, or the
shortfall residual does not transport at that threshold.

\subsection{The \texorpdfstring{$\alpha=1$}{alpha=1} Endpoint}

\begin{theorem}[$\alpha=1$ collapses to mean-CJE]
\label{prop:formal-alpha-one}
If $\Y\in[0,1]$, then $\CVaR_1(\ptarget)=V(\ptarget)$. Taking $t=1$ and
$h_1(\Sscore,\X)=1-f(\Sscore,\X)$, the shortfall audit moment is the negative
of the mean-CJE residual:
\[
\shortfall{1}{\Y}-h_1(\Sscore,\X)
=f(\Sscore,\X)-\Y
=-\varepsilon_f.
\]
The first audit moment is identically zero.
\end{theorem}

\begin{twocolproof}
\twocolrow{
\[
\Psi_{1,\ptarget}(t)
=t-\E_\ptarget[\shortfall{t}{\Y}]
=\E_\ptarget[\min\{t,\Y\}].
\]
}{
For $\alpha=1$, $t-(t-\Y)_+=\min\{t,\Y\}$ pointwise.
}
\twocolrow{
\[
\E_\ptarget[\min\{t,\Y\}]\le \E_\ptarget[\Y],
\]
with equality for any $t\ge1$.
}{
The oracle outcome is bounded in $[0,1]$.
}
\twocolrow{
\[
\CVaR_1(\ptarget)=\E_\ptarget[\Y]=V(\ptarget),
\qquad t=1\text{ is an optimizer.}
\]
}{
Take the supremum over $t$ in the previous row.
}
\twocolrow{
\[
h_1(\Sscore,\X)
=\E_{\pzero}[1-\Y\mid\mathcal{G}]
=1-f(\Sscore,\X).
\]
}{
At $t=1$, $\shortfall{1}{\Y}=1-\Y$ and conditional expectation is linear.
}
\twocolrow{
\[
g_2=(1-\Y)-(1-f(\Sscore,\X))
=f(\Sscore,\X)-\Y.
\]
}{
Substitute the $t=1$ shortfall surrogate into the audit moment.
}
\twocolrow{
\[
g_2=-(\Y-f(\Sscore,\X)),\qquad
g_1=\one\{\Y\le1\}-1=0.
\]
}{
The second moment is the negative mean-CJE residual; the first is identically
zero because $\Y\in[0,1]$.
}
\end{twocolproof}

\subsection{Discrete Outcomes and Tie Handling}
\label{app:discrete}

For atomic outcomes, the event $\{\Y\le q_\alpha\}$ can have probability
larger than $\alpha$. The conditional mean over this whole event is then not
the lower-tail expected shortfall because it averages too much mass at the
boundary.

Let
\[
q_\alpha=\inf\{t:F(t)\ge\alpha\},
\qquad
F(q_\alpha-)=\Pp(\Y<q_\alpha).
\]
The lower-tail expected shortfall with partial mass at the boundary is
\begin{equation}
\ES_{\alpha}
=\frac{1}{\alpha}\left(
\E[\Y\one\{\Y<q_\alpha\}]
+q_\alpha\left(\alpha-\Pp(\Y<q_\alpha)\right)
\right).
\label{eq:discrete-es}
\end{equation}
The plus sign before the boundary term is essential: the worst
$\alpha$ probability mass includes all outcomes below $q_\alpha$ plus only the
needed fraction of the atom at $q_\alpha$.

\begin{proposition}[Discrete stop-loss value]
\label{prop:formal-discrete-stoploss}
If $q_\alpha$ satisfies
\[
\Pp(\Y<q_\alpha)\le\alpha\le\Pp(\Y\le q_\alpha),
\]
then the optimized stop-loss value equals \cref{eq:discrete-es}.
\end{proposition}

\begin{twocolproof}
\twocolrow{
The subgradient condition above gives
\[
\Pp(\Y<q_\alpha)\le\alpha\le\Pp(\Y\le q_\alpha).
\]
}{
This is exactly the first-order condition for $q_\alpha$ to maximize the
concave stop-loss objective when the law may have an atom at $q_\alpha$.
}
\twocolrow{
Because $(q_\alpha-\Y)_+$ is zero on the boundary atom,
\[
\E[(q_\alpha-\Y)_+]
=\E[(q_\alpha-\Y)\one\{\Y<q_\alpha\}].
\]
}{
Only outcomes strictly below the threshold contribute to the lower-tail
shortfall.
}
\twocolrow{
Therefore
\[
\Psi_\alpha(q_\alpha)
=q_\alpha-\frac{1}{\alpha}
\E[(q_\alpha-\Y)\one\{\Y<q_\alpha\}].
\]
}{
Substitute the previous identity into the stop-loss objective.
}
\twocolrow{
Expanding the expectation,
\[
\Psi_\alpha(q_\alpha)
=q_\alpha-\frac{1}{\alpha}
\left(q_\alpha\Pp(\Y<q_\alpha)-\E[\Y\one\{\Y<q_\alpha\}]\right).
\]
}{
Linearity of expectation separates the threshold term from the lower-tail
mean term.
}
\twocolrow{
Rearranging,
\[
\Psi_\alpha(q_\alpha)
=\frac{1}{\alpha}
\left(\E[\Y\one\{\Y<q_\alpha\}]
+q_\alpha(\alpha-\Pp(\Y<q_\alpha))\right).
\]
}{
This is the atom-split expected shortfall formula in \cref{eq:discrete-es}.
}
\end{twocolproof}

The audit moments need a tie rule in the presence of boundary atoms. The
indicator moment with $\one\{\Y\le q_\alpha\}$ has expectation
$F(q_\alpha)-\alpha$, which need not be zero. Two valid approaches are:
\begin{itemize}[leftmargin=*,nosep]
\item use the interval condition
$\Pp(\Y<q_\alpha)\le\alpha\le\Pp(\Y\le q_\alpha)$ for the threshold moment;
\item use a randomized boundary indicator
\[
\one\{\Y<q_\alpha\}+\rho\one\{\Y=q_\alpha\}-\alpha,
\qquad
\rho=\frac{\alpha-\Pp(\Y<q_\alpha)}
{\Pp(\Y\le q_\alpha)-\Pp(\Y<q_\alpha)}.
\]
\end{itemize}
The shortfall moment remains
\[
\shortfall{q_\alpha}{\Y}-h_{q_\alpha}(\Sscore,\X),
\]
but inference must account for the chosen subgradient or tie rule. Without
that adjustment, an audit can reject solely because the outcome scale is
discrete rather than because shortfall transport failed.
